\section{A physical gap and ground-vector comparison}

The inclusive middle range has $s-6$ interactions. Distributing its lowest
energy $d$ equally therefore subtracts $d/(s-6)$ per term when the range
is nonempty; see \cite{gap:aklv13_two_block_robustness}.

\begin{definition}[A two-block truncation threshold]
    \label{def:aklv_two_block_threshold}
    \lean{TNLean.MPS.AreaLaw.aklvTwoBlockThreshold}
    \leanok
    For $g,\tau>0$ put
    \begin{align}
        \Theta(g,\tau)=\max\left\{g+11,\,
            g+8\log\frac{10^6\cdot64e^{3/4}}{\tau g}\right\}.
        \label{eq:aklv_cluster_threshold}
    \end{align}
\end{definition}

\begin{lemma}[Gap preservation and a phased ground vector for the actual truncation]
    \label{thm:aklv_two_block_scaled_low_cluster}
    \lean{TNLean.MPS.AreaLaw.exists_aklv_two_block_scaled_low_cluster}
    \leanok
    \uses{def:aklv_two_block_threshold,def:aklv_normalization_regions,
        def:aklv_two_block_truncation}
    Let $q\ge1$, $s\ge6$, $a+s+1\le L$, and let $0\le h_i\le\Id$
    be supported on $\{i,i+1\}$. Suppose $\Omega$ is a unit eigenvector of
    $H'=\sum_i h_i$ with energy $\mu$, and suppose $g>0$ satisfies
    \begin{align}
        H'-\mu\Id\ge g(\Id-|\Omega\rangle\langle\Omega|).
        \label{eq:aklv_original_rank_one_gap}
    \end{align}
    The three genuine reference minima $e_j$ and a real number $E_0$ can be
    chosen with $E_0+\sum_j e_j=\mu$ and $0\le E_0\le6$.
    They are fixed before $\tau,t$. For $0<\tau\le1$ and
    $t\ge\Theta(g,\tau)$, form the actual two-exterior family $B_T$
    with cuts $a+2,a+s-1$ and those same minima.
    Write $(v_j)$ for an orthonormal eigenbasis of $B_{\{0,1\}}$,
    with eigenvalues $\lambda_j$, and put
    $U=\operatorname{span}\{v_j:\lambda_j\le E_0+g/20\}$.
    Then $\dim U=1$, and every cutoff in
    $[E_0+g/20,E_0+19g/20]$ gives this same subspace.
    Both squared projection losses between $U$ and $\mathbb C\Omega$ are
    at most $\tau/200$.

    There is an index $j_0$ with $\lambda_{j_0}\le E_0+g/20$ such that
    every other eigenvalue satisfies
    $\lambda_j\ge\lambda_{j_0}+9g/10$.
    Thus $v_{j_0}$ is an actual unit ground vector and the truncated gap is
    at least $9g/10$. There is a real phase $\theta$ satisfying
    \begin{align}
        \|v_{j_0}-e^{i\theta}\Omega\|^2\le\frac{\tau}{100}.
        \label{eq:aklv_truncated_vector_error}
    \end{align}
    Every eigenvalue also lies in
    $(-\infty,E_0+g/10]\cup[E_0+9g/10,\infty)$.
\end{lemma}
\begin{proof}
    \leanok
    \uses{thm:aklv_normalized_two_block_residual,def:aklv_two_block_threshold}
    The three original reference minima give the centered Hamiltonian
    $H=H'-\sum_j e_j\Id$. Its constructed ground energy equals
    $E_0=\mu-\sum_j e_j$, since both energies are attained lower bounds.
    The original global gap therefore applies to $H$ at $E_0$.
    Set $\xi=64e^{3/4}e^{-(t-g)/8}$.
    The threshold (\ref{eq:aklv_cluster_threshold}) gives
    $t\ge g+11$ and $\xi\le\tau g/10^6$.
    Lemma~\ref{thm:aklv_normalized_two_block_residual} bounds the actual
    residual on each partially clipped space below $E_0+g$ by $\xi$.
    Restrict further to the original cutoff $E_0+g/40$ and the truncated
    cutoff $E_0+19g/20$.

    Using the corrected finite-window count of
    \cite{gap:alvv16_hard_truncation_windows}, the scaled low-cluster comparison
    gives equality throughout the stated
    cutoff band, preservation of its dimension, and projection losses
    $\tau/200$ in both directions. The original gap identifies its original
    low subspace with $\mathbb C\Omega$, so the new subspace has dimension one.
    Choose an eigenvector $v_{j_0}$ in it. It spans $U$; orthogonality forces
    every other eigenvector above $E_0+19g/20$. Together with
    $\lambda_{j_0}\le E_0+g/20$, this gives the gap $9g/10$.
    The projection estimate gives
    $|\langle\Omega,v_{j_0}\rangle|^2\ge1-\tau/200$.
    Choosing a minimizing relative phase yields
    (\ref{eq:aklv_truncated_vector_error}).
\end{proof}

% This is a quantitative physical consequence towards AKLV Theorem 6.1.
% It starts from original nearest-neighbour interactions, not arbitrary exterior
% operators. The source's middle shift is corrected at aggregate level, as recorded
% in docs/paper-gaps/aklv13_two_block_robustness.tex. Comparison with the printed
% sufficient threshold and elimination of tau in favor of a single t-error bound
% remain separate; the unrestricted source label is not marked complete here.
